\chapter{Chest Tightness and Discomfort}

\section*{Definition}
A subjective sensation of pressure, squeezing, constriction, or fullness in the chest area, distinct from sharp or stabbing chest pain but equally requiring careful evaluation.

\section*{What Happens in Your Body}
Chest tightness may reflect bronchoconstriction (asthma, COPD), esophageal spasm (GERD, motility disorder), anxiety/hyperventilation (muscle tension), or cardiac ischemia (angina presenting as tightness without frank pain). Differentiating cardiac from non-cardiac causes is critical.

\section*{Causes}
\begin{itemize}
  \item Asthma or reactive airway disease
  \item Gastroesophageal reflux disease (GERD) with esophageal spasm
  \item Anxiety, panic attacks, or hyperventilation syndrome
  \item Angina pectoris or coronary artery disease
  \item Chronic obstructive pulmonary disease (COPD) exacerbation
  \item Costochondritis (chest wall tenderness)
  \item Pulmonary embolism
  \item Pneumonia or pleurisy
  \item Laryngeal or tracheal irritation
  \item Heart failure (especially with orthopnea)
\end{itemize}

\section*{When to See a Doctor}
See your doctor if:
\begin{itemize}
  \item Symptoms are severe or getting worse over time
  \item Chest tightness with shortness of breath, radiating pain, sweating, nausea, or lightheadedness (possible cardiac cause — emergency)
  \item You have other concerning symptoms such as fever, weight loss, or bleeding
\end{itemize}

\section*{Related Symptoms}
\begin{itemize}
  \item Shortness of breath
  \item Wheezing
  \item Cough
  \item Anxiety
  \item Palpitations
\end{itemize}
\textbf{Important:} If this symptom persists, your doctor will check for serious underlying causes.

\section*{Self-Care / Home Management}
\begin{itemize}
  \item If chest tightness is sudden, severe, or accompanied by shortness of breath, radiating pain, sweating, or nausea, call emergency services immediately.
  \item For known asthma or COPD, use your reliever inhaler as prescribed (blue inhaler) and follow your personal asthma action plan.
  \item Sit upright, loosen tight clothing, and practise slow, deep breathing to reduce anxiety-driven chest tightness.
  \item Avoid strenuous activity and identify triggers such as cold air, allergens, or stress that provoke the sensation.
  \item Do not ignore recurrent chest tightness — keep a symptom diary of frequency, triggers, and associated symptoms to share with your doctor.
\end{itemize}

\section*{How Your Doctor Will Evaluate You}
\begin{itemize}
  \item Your doctor will differentiate cardiac from non-cardiac causes by assessing quality, duration, triggers, relieving factors, and associated symptoms.
  \item Electrocardiogram (ECG) is performed to evaluate for cardiac ischaemia, arrhythmia, or pericarditis; serial troponin levels are measured if ACS is suspected.
  \item Chest X-ray assesses for pneumonia, pneumothorax, pleural effusion, or pulmonary oedema as causes of chest tightness.
  \item Pulmonary function tests (spirometry) are arranged if asthma or COPD is suspected, with methacholine challenge for borderline cases.
\end{itemize}

\section*{Treatment Options}
\begin{itemize}
  \item Inhaled bronchodilators (salbutamol, ipratropium) and inhaled corticosteroids are first-line for asthma and COPD-related chest tightness.
  \item Anti-anginal therapy (nitrates, beta-blockers, calcium channel blockers) and antiplatelet agents (aspirin) are initiated for stable angina; PCI or CABG for acute coronary syndrome.
  \item Anxiolytic therapy (SSRIs, CBT) and breathing retraining are effective for anxiety or panic disorder presenting with chest tightness.
  \item Proton pump inhibitors or H2 receptor antagonists are prescribed for GERD-related chest tightness, often requiring an 8-week trial for symptom improvement.
  \item Medications, lifestyle modifications, and physical therapies may be prescribed as appropriate.
  \item Follow-up ensures treatment effectiveness and allows timely adjustments to the care plan.
\end{itemize}

\section*{References}
\begin{thebibliography}{99}
\bibitem{chest_tightness_and_discomfort:ref1} Global Initiative for Asthma (GINA). ``Global Strategy for Asthma Management and Prevention.'' 2024.
\bibitem{chest_tightness_and_discomfort:ref2} Global Initiative for Chronic Obstructive Lung Disease (GOLD). ``Global Strategy for the Diagnosis, Management, and Prevention of COPD.'' 2024.
\bibitem{chest_tightness_and_discomfort:ref3} Manning HL, Schwartzstein RM. ``Pathophysiology of dyspnea.'' N Engl J Med. 1995;333(23):1547-1553.
\bibitem{chest_tightness_and_discomfort:ref4} Parshall MB, Schwartzstein RM, Adams L, et al. ``An official American Thoracic Society statement: update on the mechanisms, assessment, and management of dyspnea.'' Am J Respir Crit Care Med. 2012;185(4):435-452.
\bibitem{chest_tightness_and_discomfort:ref5} Mahler DA, O'Donnell DE. ``Dyspnea: mechanisms, measurement, and management.'' 3rd ed. CRC Press; 2014.
\bibitem{chest_tightness_and_discomfort:ref6} Schwartzstein RM. ``Approach to the patient with dyspnea.'' UpToDate. Wolters Kluwer; 2025.
\end{thebibliography}
